\section{Extending the Trakr Locomotion Task to Rough Terrain}

After the flat Trakr velocity-tracking task was implemented, the existing task was extended into a rough
terrain variant. The rough task was kept separate from the flat task so that flat-ground locomotion and
terrain locomotion could be trained, evaluated, and exported independently. This was done by adding a new
robot task folder,
\texttt{trakr\_rl/tasks/locomotion/robots/trakr-rough}, while reusing the same Trakr articulation,
manager-based environment structure, PPO runner, observation design, and MDP utilities from the flat
locomotion task.

\subsection{Separate Rough-Terrain Task Registration}

The rough-terrain task was registered as a new Gymnasium environment with the id
\texttt{Trakr-Velocity-Rough}. This allowed the rough environment to appear as a separate task in the
training script while still using Isaac Lab's \texttt{ManagerBasedRLEnv}. The registration points to the
rough task's own \texttt{velocity\_env\_cfg.py}, but it continues to use the same RSL-RL PPO runner
configuration as the flat task.

\begin{lstlisting}[language=Python,caption={Rough-terrain task registration.}]
import gymnasium as gym

gym.register(
    id="Trakr-Velocity-Rough",
    entry_point="isaaclab.envs:ManagerBasedRLEnv",
    disable_env_checker=True,
    kwargs={
        "env_cfg_entry_point": f"{__name__}.velocity_env_cfg:RobotEnvCfg",
        "play_env_cfg_entry_point": f"{__name__}.velocity_env_cfg:RobotPlayEnvCfg",
        "rsl_rl_cfg_entry_point": (
            "trakr_rl.tasks.locomotion.agents.rsl_rl_ppo_cfg:"
            "BasePPORunnerCfg"
        ),
    },
)
\end{lstlisting}

With this registration, the rough-terrain policy could be trained through the same wrapper script:

\begin{lstlisting}[language=bash,caption={Training command for the rough-terrain task.}]
./trakr_rl.sh --train --task Trakr-Velocity-Rough
\end{lstlisting}

\subsection{Replacing the Flat Terrain With a Mixed Terrain Generator}

The main extension from the flat task was the terrain generator. The flat task used the same Isaac Lab
terrain importer interface, but only enabled a plane-like terrain. For rough-terrain training, the
sub-terrain dictionary was expanded to include multiple terrain types: flat ground, random height-field
roughness, pyramid slopes, ascending stairs, and descending stairs. This gave the Trakr policy exposure to
different contact patterns and body-pose disturbances during training.

\begin{lstlisting}[language=Python,caption={Mixed rough-terrain generator used by the Trakr rough task.}]
COBBLESTONE_ROAD_CFG = terrain_gen.TerrainGeneratorCfg(
    size=(8.0, 8.0),
    border_width=20.0,
    num_rows=10,
    num_cols=20,
    horizontal_scale=0.1,
    vertical_scale=0.005,
    slope_threshold=0.75,
    difficulty_range=(0.0, 1.0),
    use_cache=False,
    sub_terrains={
        "flat": terrain_gen.MeshPlaneTerrainCfg(proportion=0.1),
        "random_rough": terrain_gen.HfRandomUniformTerrainCfg(
            proportion=0.1,
            noise_range=(0.05, 0.12),
            noise_step=0.01,
            border_width=0.25,
        ),
        "hf_pyramid_slope": terrain_gen.HfPyramidSlopedTerrainCfg(
            proportion=0.1,
            slope_range=(0.5, 0.6),
            platform_width=2.0,
            border_width=0.25,
        ),
        "pyramid_stairs": terrain_gen.MeshPyramidStairsTerrainCfg(
            proportion=0.2,
            step_height_range=(0.08, 0.15),
            step_width=0.3,
            platform_width=3.0,
            border_width=1.0,
            holes=False,
        ),
        "pyramid_stairs_inv": terrain_gen.MeshInvertedPyramidStairsTerrainCfg(
            proportion=0.2,
            step_height_range=(0.08, 0.15),
            step_width=0.3,
            platform_width=3.0,
            border_width=1.0,
            holes=False,
        ),
    },
)
\end{lstlisting}

The rough task therefore preserved some flat tiles, but introduced rough height variations in the range
\texttt{0.05--0.12 m}, slopes in the range \texttt{0.5--0.6}, and stair heights in the range
\texttt{0.08--0.15 m}. The generator was arranged as a grid with \texttt{10} rows and \texttt{20} columns,
which created many terrain patches with different difficulty levels. This made it possible to use Isaac
Lab's terrain curriculum mechanism to gradually expose the policy to harder patches as performance
improved.

\subsection{Connecting the Terrain Generator to the Scene}

The rough terrain was connected to the environment through \texttt{TerrainImporterCfg}. The terrain type
was set to \texttt{generator}, and the terrain generator was assigned to the rough-terrain configuration.
The same Trakr robot asset from the flat task was then spawned on top of this generated terrain.

\begin{lstlisting}[language=Python,caption={Scene terrain importer for generated rough terrain.}]
terrain = TerrainImporterCfg(
    prim_path="/World/ground",
    terrain_type="generator",
    terrain_generator=COBBLESTONE_ROAD_CFG,
    max_init_terrain_level=1,
    collision_group=-1,
    physics_material=sim_utils.RigidBodyMaterialCfg(
        friction_combine_mode="multiply",
        restitution_combine_mode="multiply",
        static_friction=1.0,
        dynamic_friction=1.0,
    ),
    debug_vis=False,
)

robot: ArticulationCfg = ROBOT_CFG.replace(
    prim_path="{ENV_REGEX_NS}/trakr"
)
\end{lstlisting}

The height scanner and contact sensors were kept from the Trakr task and pointed to the Trakr USD prim
paths. The scanner was mounted relative to \texttt{base\_link}, while the contact sensor matched all Trakr
bodies under the robot prim. This preserved the same robot-specific body naming fixes from the flat task
while changing only the terrain distribution.

\subsection{Using Terrain Curriculum}

The rough task enabled terrain curriculum through the \texttt{terrain\_levels\_vel} curriculum term. In
the environment post-initialization step, the terrain generator curriculum flag was enabled whenever the
terrain curriculum term was present. This allowed Isaac Lab to control the terrain level assignment during
training.

\begin{lstlisting}[language=Python,caption={Enabling terrain curriculum for the rough generator.}]
class CurriculumCfg:
    terrain_levels = CurrTerm(func=mdp.terrain_levels_vel)
    lin_vel_cmd_levels = CurrTerm(mdp.lin_vel_cmd_levels)

def __post_init__(self):
    self.decimation = 4
    self.episode_length_s = 20.0
    self.sim.dt = 0.005

    if getattr(self.curriculum, "terrain_levels", None) is not None:
        if self.scene.terrain.terrain_generator is not None:
            self.scene.terrain.terrain_generator.curriculum = True
    else:
        if self.scene.terrain.terrain_generator is not None:
            self.scene.terrain.terrain_generator.curriculum = False
\end{lstlisting}

The rough environment also increased the number of parallel training environments from the flat task's
\texttt{512} environments to \texttt{2048} environments. This produced more terrain interactions per PPO
iteration, which is useful because rough-terrain locomotion has a larger state distribution than
flat-ground walking.

\begin{lstlisting}[language=Python,caption={Increased parallel environments for rough-terrain training.}]
class RobotEnvCfg(ManagerBasedRLEnvCfg):
    scene: RobotSceneCfg = RobotSceneCfg(
        num_envs=2048,
        env_spacing=2.5,
    )
\end{lstlisting}

\subsection{Command and Reward Changes for Rough Terrain}

The command distribution was narrowed for the rough-terrain experiment. Instead of training the robot on
all planar velocity directions and yaw rates, the rough task commanded lateral motion only, with
\texttt{lin\_vel\_y} between \texttt{0.5} and \texttt{1.0 m/s}. Forward velocity and yaw velocity were set
to zero. This focused training on a specific rough-terrain traversal behavior rather than a fully general
omnidirectional policy.

\begin{lstlisting}[language=Python,caption={Velocity command range used in the rough task.}]
base_velocity = mdp.UniformLevelVelocityCommandCfg(
    asset_name="robot",
    resampling_time_range=(10.0, 10.0),
    rel_standing_envs=0.1,
    debug_vis=True,
    ranges=mdp.UniformLevelVelocityCommandCfg.Ranges(
        lin_vel_x=(0.0, 0.0),
        lin_vel_y=(0.5, 1.0),
        ang_vel_z=(0.0, 0.0),
    ),
    limit_ranges=mdp.UniformLevelVelocityCommandCfg.Ranges(
        lin_vel_x=(0.0, 0.0),
        lin_vel_y=(0.5, 1.0),
        ang_vel_z=(0.0, 0.0),
    ),
)
\end{lstlisting}

The reward weights were also adjusted for rough terrain. The linear velocity tracking reward was increased
to encourage stronger command following, while action-rate, roll, pitch-rate, foot air-time variance, foot
height, and undesired-contact penalties were used to discourage unstable terrain interactions. A
body-frame foot-height penalty was added so that the legs maintained useful clearance over steps and
uneven surfaces.

\begin{lstlisting}[language=Python,caption={Selected rough-terrain reward terms.}]
track_lin_vel_xy = RewTerm(
    func=mdp.track_lin_vel_xy_exp,
    weight=4.0,
    params={"command_name": "base_velocity", "std": math.sqrt(0.25)},
)

roll_penalty = RewTerm(func=mdp.roll_penalty, weight=-1.0)
pitch_penalty = RewTerm(func=mdp.pitch_rate_penalty, weight=-1.0)

air_time_variance = RewTerm(
    func=mdp.air_time_variance_penalty,
    weight=-0.75,
    params={"sensor_cfg": SceneEntityCfg("contact_forces", body_names=".*_toe")},
)

feet_height_body = RewTerm(
    func=mdp.feet_height_body,
    weight=-2.0,
    params={
        "command_name": "base_velocity",
        "target_height": 0.15,
        "tanh_mult": 8.0,
        "asset_cfg": SceneEntityCfg("robot", body_names=[".*_toe"]),
    },
)
\end{lstlisting}

\subsection{Play Configuration for Terrain Inspection}

The play configuration was reduced to a small number of environments so the trained policy could be
inspected visually. During play, the environment uses \texttt{4} robots and a \texttt{4 \(\times\) 4}
terrain layout, and the command ranges are set directly to the final limit ranges. This makes it easier to
observe the robot crossing slopes, rough patches, and stair terrain without launching the full training
batch.

\begin{lstlisting}[language=Python,caption={Reduced play configuration for rough-terrain inspection.}]
class RobotPlayEnvCfg(RobotEnvCfg):
    def __post_init__(self):
        super().__post_init__()
        self.scene.num_envs = 4
        self.scene.terrain.terrain_generator.num_rows = 4
        self.scene.terrain.terrain_generator.num_cols = 4
        self.commands.base_velocity.ranges = (
            self.commands.base_velocity.limit_ranges
        )
\end{lstlisting}

Overall, the rough-terrain extension was implemented by taking the already working Trakr flat locomotion
task and replacing the flat-only terrain distribution with a mixed generated terrain curriculum. The robot
asset, observation structure, contact sensors, action space, and PPO runner remained largely the same,
while the terrain generator, command range, reward weights, environment count, and play layout were
modified to train the Trakr policy on random roughness, slopes, ascending stairs, and descending stairs.
